% !TEX root = ../main.tex

\section{Task A: Read Data from the LIS3MDL Magnetic Sensor}

\subsection{a) Calculations:}

In the implemented startup sequence, the sensor is configured by writing values to \texttt{CTRL\_REG1} to \texttt{CTRL\_REG4} via I\textsuperscript{2}C. A \texttt{WHO\_AM\_I} check would be useful as an additional verification step, but it is not used in the present implementation. The relevant registers are listed below.

\begin{table}[H]
	\centering
	\begin{tabularx}{\textwidth}{|>{\raggedright\arraybackslash}X|>{\raggedright\arraybackslash}p{3cm}|>{\raggedright\arraybackslash}p{2.4cm}|>{\raggedright\arraybackslash}p{4.2cm}|}
		\hline
		\textbf{Parameter Description}                                            & \textbf{Value}          & \textbf{Unit (if any)} & \textbf{Register Address} \\ \hline
		Read device identity (\texttt{WHO\_AM\_I}), expected response             & optional, \texttt{0x3D} & --                     & \texttt{0x0F}             \\ \hline
		Configure X/Y performance mode and output data rate (\texttt{CTRL\_REG1}) & \texttt{0xFC}           & --                     & \texttt{0x20}             \\ \hline
		Configure full-scale range (\texttt{CTRL\_REG2})                          & \texttt{0x40}           & $\pm 12$~Gauss         & \texttt{0x21}             \\ \hline
		Enable continuous-conversion mode (\texttt{CTRL\_REG3})                   & \texttt{0x00}           & --                     & \texttt{0x22}             \\ \hline
		Configure Z-axis performance mode (\texttt{CTRL\_REG4})                   & \texttt{0x0C}           & --                     & \texttt{0x23}             \\ \hline
		Optional: block data update / endianness (\texttt{CTRL\_REG5})            & not used                & --                     & \texttt{0x24}             \\ \hline
	\end{tabularx}
\end{table}

The meaning of the configured commands is:
\begin{itemize}
	\item \texttt{CTRL\_REG1} sets temperature compensation, X/Y operating mode, and output data rate.
	\item \texttt{CTRL\_REG2} selects the full-scale range. In this implementation, \texttt{0x40} selects \(\pm 12\)~Gauss.
	\item \texttt{CTRL\_REG3} selects the operating mode. \texttt{0x00} means continuous-conversion mode.
	\item \texttt{CTRL\_REG4} sets the Z-axis operating mode.
\end{itemize}

\begin{table}[H]
	\centering
	\begin{tabularx}{0.85\textwidth}{|>{\raggedright\arraybackslash}X|>{\raggedright\arraybackslash}X|}
		\hline
		\textbf{Description}                                            & \textbf{Register Address} \\ \hline
		X-axis low byte (\texttt{SENSOR\_OUT\_X\_L})                    & \texttt{0x28}             \\ \hline
		X-axis high byte (\texttt{SENSOR\_OUT\_X\_H})                   & \texttt{0x29}             \\ \hline
		Y-axis low byte (\texttt{SENSOR\_OUT\_Y\_L})                    & \texttt{0x2A}             \\ \hline
		Y-axis high byte (\texttt{SENSOR\_OUT\_Y\_H})                   & \texttt{0x2B}             \\ \hline
		Z-axis low byte (\texttt{SENSOR\_OUT\_Z\_L})                    & \texttt{0x2C}             \\ \hline
		Z-axis high byte (\texttt{SENSOR\_OUT\_Z\_H})                   & \texttt{0x2D}             \\ \hline
		Status register (\texttt{STATUS\_REG}, optional before readout) & \texttt{0x27}             \\ \hline
	\end{tabularx}
\end{table}

Each axis value is stored in two 8-bit registers. These two bytes are combined into one signed 16-bit value in two's complement representation:
\[
	M_{x,\mathrm{raw}} = (\texttt{OUT\_X\_H} \ll 8) \; | \; \texttt{OUT\_X\_L}
\]
\[
	M_{y,\mathrm{raw}} = (\texttt{OUT\_Y\_H} \ll 8) \; | \; \texttt{OUT\_Y\_L}
\]
\[
	M_{z,\mathrm{raw}} = (\texttt{OUT\_Z\_H} \ll 8) \; | \; \texttt{OUT\_Z\_L}
\]

After combining, the result must be interpreted as a signed 16-bit integer:
\[
	M_{\mathrm{raw}} \in [-32768,\,32767]
\]

For the configured full-scale range \(\pm 12\)~Gauss, the sensitivity of the LIS3MDL is
\[
	S = 2281~\mathrm{LSB/Gauss}.
\]

Therefore, the magnetic field value is
\[
	M_{\mathrm{Gauss}} = \frac{M_{\mathrm{raw}}}{2281}.
\]

In the program, the value is first converted to milli-Gauss in order to print three decimal places without floating-point arithmetic:
\[
	M_{\mathrm{mGauss}} = \frac{M_{\mathrm{raw}} \cdot 1000}{2281}.
\]

Thus, for each axis:
\[
	M_{x,\mathrm{mGauss}} = \frac{M_{x,\mathrm{raw}} \cdot 1000}{2281}, \qquad
	M_{y,\mathrm{mGauss}} = \frac{M_{y,\mathrm{raw}} \cdot 1000}{2281}, \qquad
	M_{z,\mathrm{mGauss}} = \frac{M_{z,\mathrm{raw}} \cdot 1000}{2281}.
\]

The value is then split into integer part and fractional part:
\[
	\texttt{intPart} = \left\lfloor \frac{M_{\mathrm{mGauss}}}{1000} \right\rfloor,
	\qquad
	\texttt{fracPart} = M_{\mathrm{mGauss}} \bmod 1000.
\]

The implemented error handling checks whether the converted value exceeds the selected full-scale range:
\[
	M_{\mathrm{mGauss}} \notin [-12000,\;12000]
\]
which corresponds to
\[
	M_{\mathrm{Gauss}} \notin [-12,\;12].
\]

Therefore, the program must:
\begin{enumerate}
	\item write the startup configuration to \texttt{CTRL\_REG1} to \texttt{CTRL\_REG4},
	\item read two bytes per axis via I\textsuperscript{2}C,
	\item combine them to one signed 16-bit value,
	\item convert the raw value using the sensitivity \(2281~\mathrm{LSB/Gauss}\),
	\item print the result via UART,
	\item report an error if the measured value exceeds \(\pm 12\)~Gauss.
\end{enumerate}

\vspace{1.5cm}

\subsection{b) Implementation:}

At first, UART4 and I2C2 were configured in STM32CubeMX. UART4 is used to transmit the measurement results to the serial monitor, and I2C2 is used to communicate with the LIS3MDL sensor.

\begin{figure}[H]
	\begin{center}
		\includegraphics[width=0.75\textwidth]{TaskA_URAT_Conf.png}
	\end{center}
	\caption{CubeMX configuration for UART4}
\end{figure}

\begin{figure}[H]
	\begin{center}
		\includegraphics[width=0.75\textwidth]{TaskA_I2C_Conf1.png}
	\end{center}
	\caption{CubeMX configuration for I2C2}
\end{figure}

\begin{figure}[H]
	\begin{center}
		\includegraphics[width=0.9\textwidth]{TaskA_I2C_Conf2.png}
	\end{center}
	\caption{I2C GPIO configuration}
\end{figure}

The code was inserted into several STM32CubeMX user blocks. The corresponding block names are kept in the listings.

\paragraph{Additional include files (\texttt{USER CODE BEGIN Includes})}
The libraries \texttt{string.h} and \texttt{stdio.h} are needed for string formatting and UART message generation.

\begin{lstlisting}[language=C,caption={Additional include files in USER CODE BEGIN Includes}]
/* USER CODE BEGIN Includes */
#include "string.h"
#include "stdio.h"
/* USER CODE END Includes */
\end{lstlisting}

\paragraph{Register definitions (\texttt{USER CODE BEGIN PD})}
All device addresses and register addresses are defined as macros. This makes the code more readable and reusable.

\begin{lstlisting}[language=C,caption={Register and address macros in USER CODE BEGIN PD}]
/* USER CODE BEGIN PD */
// Sensor device addresses
#define SENSOR_WRITE_ADDR   0x3C
#define SENSOR_READ_ADDR    0x3D

// Control registers
#define SENSOR_CTRL_REG1    0x20
#define SENSOR_CTRL_REG2    0x21
#define SENSOR_CTRL_REG3    0x22
#define SENSOR_CTRL_REG4    0x23

// Output registers for raw magnetic field data
#define SENSOR_OUT_X_L      0x28
#define SENSOR_OUT_X_H      0x29
#define SENSOR_OUT_Y_L      0x2A
#define SENSOR_OUT_Y_H      0x2B
#define SENSOR_OUT_Z_L      0x2C
#define SENSOR_OUT_Z_H      0x2D
/* USER CODE END PD */
\end{lstlisting}

\paragraph{Function prototypes (\texttt{USER CODE BEGIN PFP})}
Two function prototypes are declared: one generic readout function for one axis and one basic plausibility check for the converted field value.

\begin{lstlisting}[language=C,caption={Function prototypes in USER CODE BEGIN PFP}]
/* USER CODE BEGIN PFP */
void getMagField(uint16_t addressAxisL, uint16_t addressAxisH,
                 uint16_t divFactor, char axis);
void checkForImpossibleVal(int32_t magneticField_mGauss);
/* USER CODE END PFP */
\end{lstlisting}

\paragraph{Generic measurement and error-check functions (\texttt{USER CODE BEGIN 0})}
The function \texttt{getMagField()} reads the low byte and high byte of one axis, combines them to a signed 16-bit raw value, converts the value to milli-Gauss, checks whether the value is plausible, formats the output string, and sends it via UART4. The function \texttt{checkForImpossibleVal()} implements the basic error handling by checking whether the value exceeds the configured range of \(\pm 12\)~Gauss.

\begin{lstlisting}[language=C,caption={Measurement and basic error handling in USER CODE BEGIN 0}]
/* USER CODE BEGIN 0 */
void getMagField(uint16_t addressAxisL, uint16_t addressAxisH, uint16_t divFactor, char axis)
{
    uint8_t outLsb = 0;
    uint8_t outMsb = 0;
    int16_t rawOutput = 0;
    int32_t magneticField_mGauss = 0;
    int32_t intPart = 0;
    int32_t fracPart = 0;
    char msg[40] = {'\0'};

    HAL_I2C_Mem_Read(&hi2c2, SENSOR_READ_ADDR, addressAxisL,
                     I2C_MEMADD_SIZE_8BIT, &outLsb, 1, HAL_MAX_DELAY);
    HAL_I2C_Mem_Read(&hi2c2, SENSOR_READ_ADDR, addressAxisH,
                     I2C_MEMADD_SIZE_8BIT, &outMsb, 1, HAL_MAX_DELAY);

    rawOutput = (int16_t)(((uint16_t)outMsb << 8) | outLsb);

    magneticField_mGauss = ((int32_t)rawOutput * 1000) / divFactor;
    checkForImpossibleVal(magneticField_mGauss);
    intPart = magneticField_mGauss / 1000;
    fracPart = magneticField_mGauss % 1000;

    if (fracPart < 0)
    {
        fracPart = -fracPart;
    }

    if (fracPart < 0)
        fracPart = -fracPart;

    if ((magneticField_mGauss < 0) && (intPart == 0))
    {
        if (axis == 'Z')
            sprintf(msg, "%c: -0.%03d Gauss\r\n", axis, (int)fracPart);
        else
            sprintf(msg, "%c: -0.%03d Gauss, ", axis, (int)fracPart);
    }
    else
    {
        if (axis == 'Z')
            sprintf(msg, "%c: %d.%03d Gauss\r\n", axis, (int)intPart, (int)fracPart);
        else
            sprintf(msg, "%c: %d.%03d Gauss, ", axis, (int)intPart, (int)fracPart);
    }

    HAL_UART_Transmit(&huart4, (uint8_t *)msg, strlen(msg), HAL_MAX_DELAY);

    HAL_Delay(200);
}

void checkForImpossibleVal(int32_t magneticField_mGauss)
{
    if (magneticField_mGauss > 12000 || magneticField_mGauss < -12000)
    {
        char msg[64] = {'\0'};
        sprintf(msg, "Error magnetic field value > 12 Gauss\r\n");
        HAL_UART_Transmit(&huart4, (uint8_t *)msg, strlen(msg), HAL_MAX_DELAY);
        HAL_Delay(200);
    }
}
/* USER CODE END 0 */
\end{lstlisting}

\paragraph{Sensor startup sequence (\texttt{USER CODE BEGIN 2})}
The startup sequence is executed once after peripheral initialization. The configuration bytes are stored in variables and written to the sensor registers via \texttt{HAL\_I2C\_Mem\_Write()}.

\begin{lstlisting}[language=C,caption={Sensor initialization in USER CODE BEGIN 2}]
/* USER CODE BEGIN 2 */
uint8_t xyConfig = 0xFC;
uint8_t zConfig = 0x0C;
uint8_t scaleConfig = 0x40;
uint8_t modeConfig = 0x00;
uint16_t gain = 2281;

HAL_I2C_Mem_Write(&hi2c2, SENSOR_WRITE_ADDR, SENSOR_CTRL_REG1,
                  I2C_MEMADD_SIZE_8BIT, &xyConfig, 1, HAL_MAX_DELAY);
HAL_I2C_Mem_Write(&hi2c2, SENSOR_WRITE_ADDR, SENSOR_CTRL_REG4,
                  I2C_MEMADD_SIZE_8BIT, &zConfig, 1, HAL_MAX_DELAY);

HAL_I2C_Mem_Write(&hi2c2, SENSOR_WRITE_ADDR, SENSOR_CTRL_REG2,
                  I2C_MEMADD_SIZE_8BIT, &scaleConfig, 1, HAL_MAX_DELAY);
HAL_I2C_Mem_Write(&hi2c2, SENSOR_WRITE_ADDR, SENSOR_CTRL_REG3,
                  I2C_MEMADD_SIZE_8BIT, &modeConfig, 1, HAL_MAX_DELAY);
/* USER CODE END 2 */
\end{lstlisting}

\paragraph{Main loop}
Inside the infinite loop, the same generic function is called three times, once for each axis.

\begin{lstlisting}[language=C,caption={Three-axis readout in the while-loop}]
/* USER CODE END WHILE */
getMagField(SENSOR_OUT_X_L,SENSOR_OUT_X_H,gain,'X');
getMagField(SENSOR_OUT_Y_L,SENSOR_OUT_Y_H,gain,'Y');
getMagField(SENSOR_OUT_Z_L,SENSOR_OUT_Z_H,gain,'Z');
/* USER CODE BEGIN 3 */
\end{lstlisting}

The resulting program flow is:
\begin{enumerate}
	\item initialize the MCU peripherals generated by CubeMX,
	\item configure the LIS3MDL once in \texttt{USER CODE BEGIN 2},
	\item read the X-, Y-, and Z-axis values continuously,
	\item check whether a measured value exceeds the configured range,
	\item transmit the formatted values via UART4.
\end{enumerate}

\vspace{1.5cm}

\subsection{c) Results:}

The UART output on the serial monitor is shown below. The values change over time when the board is moved. In the shown recording, some lines are incomplete because the serial monitor was opened or closed during an ongoing transmission.

\begin{lstlisting}[caption={UART output of the magnetic field values}]
---- Opened the serial port /dev/ttyACM1 ----
: -0.448 Gauss
X: 0.535 Gauss, Y: -0.101 Gauss, Z: -0.455 Gauss
X: 0.539 Gauss, Y: -0.104 Gauss, Z: -0.455 Gauss
X: 0.544 Gauss, Y: -0.110 Gauss, Z: -0.455 Gauss
X: 0.543 Gauss, Y: -0.109 Gauss, 
---- Closed the serial port /dev/ttyACM1 ----
\end{lstlisting}

The intended output format of one complete measurement line is
\begin{lstlisting}[caption={Intended UART output format}]
X: 0.535 Gauss, Y: -0.101 Gauss, Z: -0.455 Gauss
\end{lstlisting}

\vspace{1.5cm}

\subsection{d) Discussion:}

A central design decision was to implement the readout in a generic function instead of writing separate code for the X-, Y-, and Z-axis. This reduces redundancy and improves readability. Only the register addresses, the gain factor, and the axis character have to be passed as parameters.

The basic error handling was implemented in a separate function. It checks whether the calculated magnetic field value exceeds the configured full-scale range of \(\pm 12\)~Gauss. This is a plausibility check and helps to detect obviously invalid measurement results. However, it does not check whether the I\textsuperscript{2}C communication itself failed, because the return values of \texttt{HAL\_I2C\_Mem\_Read()} and \texttt{HAL\_I2C\_Mem\_Write()} are not evaluated.

During implementation, the most critical points were:
\begin{itemize}
	\item correct register addresses for all three axes,
	\item correct byte order when combining low byte and high byte,
	\item correct interpretation as signed 16-bit value,
	\item correct gain factor for the configured range,
	\item correct handling of negative values in the UART output.
\end{itemize}

The most reusable parts of the code are:
\begin{itemize}
	\item the macro definitions for register and device addresses,
	\item the generic function \texttt{getMagField()},
	\item the plausibility check \texttt{checkForImpossibleVal()},
	\item the startup pattern using \texttt{HAL\_I2C\_Mem\_Write()},
	\item the UART output routine based on \texttt{sprintf()} and \texttt{HAL\_UART\_Transmit()}.
\end{itemize}

A possible improvement would be to add a \texttt{WHO\_AM\_I} check during startup and to evaluate the return values of the HAL I\textsuperscript{2}C functions in order to detect communication errors directly.
